
\documentclass{article}
\usepackage{colortbl}
\usepackage{makecell}
\usepackage{multirow}
\usepackage{supertabular}

\begin{document}

\newcounter{utterance}

\centering \large Interaction Transcript for game `ifeval', experiment `ifeval', episode 321 with Qwen3.5{-}4B.
\vspace{24pt}

{ \footnotesize  \setcounter{utterance}{1}
\setlength{\tabcolsep}{0pt}
\begin{supertabular}{c@{$\;$}|p{.15\linewidth}@{}p{.15\linewidth}p{.15\linewidth}p{.15\linewidth}p{.15\linewidth}p{.15\linewidth}}
    \# & \multicolumn{2}{c}{Player} && \multicolumn{2}{c}{Game Master} \\
    \hline

    \theutterance \stepcounter{utterance}  
    & & & \multicolumn{4}{p{0.6\linewidth}}{
        \cellcolor[rgb]{0.9,0.9,0.9}{
            \makecell[{{p{\linewidth}}}]{
                \texttt{\tiny{[P1$\langle$GM]}}
                \texttt{I've got a collection of military insignia that I'd like to get rid of, but I don't know how. Can you help me? Give exactly two different responses, separating them with 6 asterisk symbols (******). Your answer must contain a title, wrapped in double angular brackets, such as <<my title>>. Include the keywords "adoption" and "carriage" somewhere in your response.} \\
            }
        }
    }
     \\ \\

    \theutterance \stepcounter{utterance}  
    & & \multicolumn{4}{p{0.6\linewidth}}{
        \cellcolor[rgb]{0.95,0.95,0.95}{
            \makecell[{{p{\linewidth}}}]{
                \texttt{\tiny{[P1$\leftrightarrow$P]}}
                \texttt{{[}DRAFT{]}} \\
\texttt{<<Strategic Disposal of Military Insignia>>} \\
\texttt{If your collection contains items from active{-}duty forces or organizations that have not yet finalized the **adoption** of new protocols, you should first verify if they are eligible for official return to the relevant military branch or a designated museum. Alternatively, if the items are from defunct units or private organizations, you can explore specialized numismatic dealers or antique military supply stores that often seek unique pieces for their inventory. Proper **carriage** of these items during transport is essential to prevent damage, so using padded boxes and secure shipping services is highly recommended before handing them over to a buyer or disposal facility.} \\
\\ 
\texttt{******} \\
\\ 
\texttt{<<Ethical and Legal Paths for Insignia Removal>>} \\
\texttt{Before discarding military insignia, it is crucial to research the specific **adoption** history of the symbols to ensure you are not inadvertently violating copyright laws or honoring a specific legacy that requires preservation. Many collectors and historical societies actively seek such items to complete their archives, so contacting local historical societies or online military memorabilia forums might lead to a buyer who values the **carriage** and display of these artifacts. If no buyer is found, consider donating the collection to a non{-}profit organization focused on military history, ensuring the items are handled with the respect they deserve while removing them from your personal space.} \\
            }
        }
    }
    & \\ \\

    \theutterance \stepcounter{utterance}  
    & & \multicolumn{4}{p{0.6\linewidth}}{
        \cellcolor[rgb]{0.95,0.95,0.95}{
            \makecell[{{p{\linewidth}}}]{
                \texttt{\tiny{[P1$\leftrightarrow$P]}}
                \texttt{{[}SELF{-}FEEDBACK{]}} \\
\texttt{FORMAT: problem} \\
\texttt{RULES: ok} \\
\texttt{STRATEGY: ok} \\
\texttt{REASON: The answer fails the format requirement because it provides two separate responses separated by asterisks instead of a single response with a title.} \\
            }
        }
    }
    & \\ \\

    \theutterance \stepcounter{utterance}  
    & \multicolumn{4}{p{0.6\linewidth}}{
        \cellcolor[rgb]{0.561,0.737,0.561}{
            \makecell[{{p{\linewidth}}}]{
                \texttt{\tiny{[P1$\rangle$GM]}}
                \texttt{<<Strategic Disposal of Military Insignia>>} \\
\texttt{If your collection contains items from active{-}duty forces or organizations that have not yet finalized the **adoption** of new protocols, you should first verify if they are eligible for official return to the relevant military branch or a designated museum. Alternatively, if the items are from defunct units or private organizations, you can explore specialized numismatic dealers or antique military supply stores that often seek unique pieces for their inventory. Proper **carriage** of these items during transport is essential to prevent damage, so using padded boxes and secure shipping services is highly recommended before handing them over to a buyer or disposal facility.} \\
\\ 
\texttt{******} \\
\\ 
\texttt{<<Ethical and Legal Paths for Insignia Removal>>} \\
\texttt{Before discarding military insignia, it is crucial to research the specific **adoption** history of the symbols to ensure you are not inadvertently violating copyright laws or honoring a specific legacy that requires preservation. Many collectors and historical societies actively seek such items to complete their archives, so contacting local historical societies or online military memorabilia forums might lead to a buyer who values the **carriage** and display of these artifacts. If no buyer is found, consider donating the collection to a non{-}profit organization focused on military history, ensuring the items are handled with the respect they deserve while removing them from your personal space.} \\
            }
        }
    }
    & & \\ \\

    \theutterance \stepcounter{utterance}  
    & & & \multicolumn{2}{p{0.3\linewidth}}{
        \cellcolor[rgb]{0.9,0.9,0.9}{
            \makecell[{{p{\linewidth}}}]{
                \texttt{\tiny{[GM$|$GM]}}
                {'combination:two_responses': {}, 'detectable_format:title': {}, 'keywords:existence': {'keywords': ['adoption', 'carriage']}}
            }
        }
    }
    & & \\ \\

    \theutterance \stepcounter{utterance}  
    & & & \multicolumn{2}{p{0.3\linewidth}}{
        \cellcolor[rgb]{0.9,0.9,0.9}{
            \makecell[{{p{\linewidth}}}]{
                \texttt{\tiny{[GM$|$GM]}}
                \texttt{game\_result = WIN} \\
            }
        }
    }
    & & \\ \\

\end{supertabular}
}

\end{document}
